% fig01.tex -- Set A (Architectural) Figure 1: theorem ladder + post-core extension map.

\centering
\begin{tikzpicture}[setA]

% ============================================================
% Styles.
%   Panel A (the core ladder): chip + slab + premise note.
%   Panel B (the post-core extension map): titled cards with
%   per-item status chips.
% ============================================================
\tikzset{
  ladderChip/.style={layerChip, minimum width=11mm, minimum height=11mm},
  ladderSlab/.style={block, text width=82mm, minimum height=11mm,
                     align=left, inner sep=4pt, font=\sffamily\footnotesize},
  ladderNote/.style={font=\sffamily\scriptsize\itshape, text=gMid,
                     text width=44mm, align=left, anchor=west,
                     inner sep=0pt},
  panelHd/.style={font=\sffamily\bfseries\footnotesize, text=gDark,
                  anchor=west},
  cardTtl/.style={font=\sffamily\bfseries\footnotesize, text=gDark,
                  anchor=north west, text width=62mm, align=left,
                  inner sep=0pt},
  cardAtt/.style={font=\sffamily\scriptsize\itshape, text=gMid,
                  anchor=north west, text width=62mm, align=left,
                  inner sep=0pt},
  cardTxt/.style={font=\sffamily\scriptsize, text=gDark,
                  anchor=north west, text width=62mm, align=left,
                  inner sep=0pt},
  cardTxtG/.style={font=\sffamily\scriptsize, text=gMid,
                   anchor=north west, text width=62mm, align=left,
                   inner sep=0pt},
  chipNow/.style={draw=gDark, line width=0.5pt, rounded corners=2pt,
                  fill=gPale, font=\sffamily\bfseries\scriptsize,
                  text=gDark, inner xsep=4pt, inner ysep=1.6pt,
                  anchor=north west},
  chipFut/.style={draw=gLight, line width=0.5pt, rounded corners=2pt,
                  fill=white, font=\sffamily\bfseries\scriptsize,
                  text=gMid, inner xsep=4pt, inner ysep=1.6pt,
                  anchor=north west},
  cardBox/.style={draw=gMid, line width=0.6pt, rounded corners=2.5pt,
                  fill=gBg},
}

\def\rowGap{2.5mm}

% ============================================================
% PANEL A -- the source-side core: the theorem ladder.
% Layers 0--4, each rung Conditional and depending on the rung
% above it.
% ============================================================
\node[panelHd] (hdA) at (-5.5mm,10mm)
  {The theorem ladder \textemdash\ source-side core (Layers 0--4)};

% --- LAYER 0/1 (two formal layers folded into one row) ---
\node[ladderChip] (c01) at (0,0) {0\,/\,1};
\node[ladderSlab,  right=1.5mm of c01, anchor=west] (s01) {%
  \textbf{Source histories, replay events, abstract DBLog runs.}
  Defines $\WellformedRun$ (source-log retention, frontier discipline,
  watermark consistency, chunk-read honesty).};
\node[ladderNote,  right=2mm   of s01]              (n01)                      {%
  Conditional on the wellformed-run model.\\
  WF carries the source-side observation\\
  obligations; the certificate's external-\\
  observation premise enters at Layer 3.};

% --- LAYER 2 ---
\node[ladderChip,  below=\rowGap of c01.south, anchor=north] (c2)              {2};
\node[ladderSlab,  right=1.5mm of c2,  anchor=west]          (s2)              {%
  \textbf{Source-side virtual-cut soundness.}
  $\WellformedRun \Rightarrow$ replay at frontier equals source on the
  run scope.};
\node[ladderNote,  right=2mm   of s2]                        (n2)              {%
  Same premises as L0/1; no \emph{new}\\
  assumptions (external observation\\
  is already inherited).};

% --- LAYER 3 ---
\node[ladderChip,  below=\rowGap of c2.south, anchor=north]  (c3)              {3};
\node[ladderSlab,  right=1.5mm of c3, anchor=west]           (s3)              {%
  \textbf{Accepted-certificate soundness.}
  $\Verify(C{,}E)=\Accept \,\land\, \FaithfulSourceObservation \,\land\,$
  checker substrate $\Rightarrow \VirtualCut(C, H)$.};
\node[ladderNote,  right=2mm   of s3]                        (n3)              {%
  Conditional on Layer 2 $+$\\
  checker-substrate obligations $+$\\
  External observation assumption.};

% --- LAYER 4 ---
\node[ladderChip,  below=\rowGap of c3.south, anchor=north]  (c4)              {4};
\node[ladderSlab,  right=1.5mm of c4, anchor=west]           (s4)              {%
  \textbf{Whole-table specialization.}
  $\VirtualCut(C, H) \,\land\, \WholeTableClaimScope(C, H, b) \Rightarrow$
  unrestricted equality on the table at the frontier.};
\node[ladderNote,  right=2mm   of s4]                        (n4)              {%
  Conditional on Layer 3 $+$\\
  $\WholeTableClaimScope(C, H, b)$.};

% Dependency spine: lower depends on upper.
\foreach \a/\b in {c01/c2, c2/c3, c3/c4} {
  \draw[-{Latex[length=5pt, width=4pt]}, gDark, line width=0.9pt]
    ([yshift=-0.2pt]\a.south)
    -- ([yshift=0.2pt]\b.north);
}

% ============================================================
% Panel divider -- a neutral layout rule. NOT a status boundary:
% the post-core area below is mixed-status, so no "future work"
% line is drawn.
% ============================================================
\coordinate (divL) at ($(c4.south west) + (-6.5mm, -7mm)$);
\coordinate (divR) at ($(n4.east |- divL) + (0.5mm, 0)$);
\draw[gLight, line width=0.4pt] (divL) -- (divR);

% ============================================================
% PANEL B -- the post-core extension map.
% Two independent extension paths off the core; NOT ladder
% rungs. Each card carries per-item status chips.
% ============================================================
\node[panelHd, anchor=north west] (hdB) at ($(divL) + (0,-4mm)$)
  {Post-core extension map \textemdash\ independent extensions, \emph{not} ladder rungs};

% --- Card 1: source-side virtual-cut algebra (extends Layer 2) ---
\node[cardTtl, anchor=north west] (a-ttl) at ($(hdB.south west) + (1.5mm,-3.5mm)$)
  {Source-side virtual-cut algebra};
\node[cardAtt, anchor=north west] (a-att) at ([yshift=-0.8mm]a-ttl.south west)
  {Extends the Layer~2 virtual-cut theorem.};
\node[chipNow, anchor=north west] (a-c1) at ([yshift=-1.5mm]a-att.south west)
  {Conditional};
\node[cardTxt, anchor=north west] (a-d1) at ([yshift=-0.8mm]a-c1.south west)
  {Continuation across frontiers; restriction to a sub-scope.};
\node[chipFut, anchor=north west] (a-c2) at ([yshift=-1.6mm]a-d1.south west)
  {Future work};
\node[cardTxtG, anchor=north west] (a-d2) at ([yshift=-0.8mm]a-c2.south west)
  {Disjoint composition; concrete extended-certificate verification;
   raw overlap/retry normalization.};

% --- Card 2: destination-side sink extensions (extend the certified prefix) ---
\node[cardTtl, anchor=north west] (b-ttl) at ($(a-ttl.north west) + (72mm,0)$)
  {Destination-side sink extensions};
\node[cardAtt, anchor=north west] (b-att) at ([yshift=-0.8mm]b-ttl.south west)
  {Extend the core's certified clean prefix; need a separate sink model.};
\node[chipFut, anchor=north west] (b-c1) at ([yshift=-1.5mm]b-att.south west)
  {Future work, Non-claim};
\node[cardTxtG, anchor=north west] (b-d1) at ([yshift=-0.8mm]b-c1.south west)
  {ACK-only receipt; stream-prefix correctness. Applied destination
   state is a Non-claim.};
\node[chipFut, anchor=north west] (b-c2) at ([yshift=-1.6mm]b-d1.south west)
  {Future work};
\node[cardTxtG, anchor=north west] (b-d2) at ([yshift=-0.8mm]b-c2.south west)
  {Strong sink-state barrier convergence; needs a fenced, idempotent
   sink contract and visible apply-barrier evidence.};

% Card background boxes (drawn behind the content).
\begin{scope}[on background layer]
  \node[cardBox, fit=(a-ttl)(a-att)(a-c1)(a-d1)(a-c2)(a-d2),
        inner sep=4pt] {};
  \node[cardBox, fit=(b-ttl)(b-att)(b-c1)(b-d1)(b-c2)(b-d2),
        inner sep=4pt] {};
\end{scope}

\end{tikzpicture}

\caption{Theorem ladder and post-core extension map. \textbf{Top
panel --- the source-side core.} Layers~0--4 form the core dependency
ladder: each layer is Conditional on the premises listed beside it
and on every layer above it. \textbf{Bottom panel --- the post-core
extension map}, \emph{not} a continuation of the ladder. The
source-side virtual-cut algebra card extends the Layer~2 virtual-cut
theorem: continuation across later frontiers is Conditional on that
theorem, a wellformed source history, and the faithful
continuation-segment premise, while restriction to a sub-scope is
Conditional only on an existing cut and a sub-scope inclusion;
disjoint composition, concrete
extended-certificate verification, and raw overlap/retry
normalization remain Future work. The destination-side sink
extensions extend the certified clean prefix the core produces but
require a separate sink model: ACK-only receipt and strong sink-state
barrier convergence are Future work, and applied destination-state
correctness is a Non-claim without a fenced, idempotent sink contract
and visible apply-barrier evidence. When a virtual cut is obtained
from an accepted certificate, the accepted-certificate premises and
the external observation assumption are inherited.}
\label{fig:theorem-ladder}
